\worksheettitle
  {M04-R. Сначала порядок}
  {\modulemark{M04-R} Коррекционная карточка}

\studentnameblock

\begin{warningbox}[Назовите ошибку]
  Нельзя выбирать действие, пока не установлено, как расположены точки и
  какой отрезок является целым.
\end{warningbox}

\begin{practicebox}[Шаг 1. Прочитайте схему]
  Для порядка $A-C-B$:
  \[
    \text{целый отрезок } \answerblank[20mm],\qquad
    \text{части } \answerblank[18mm]\text{ и }\answerblank[18mm].
  \]
\end{practicebox}

\begin{practicebox}[Шаг 2. Запишите равенство]
  \[
    AB=\answerblank[18mm]+\answerblank[18mm].
  \]
  Если неизвестна часть $CB$, то
  \[
    CB=\answerblank[18mm]-\answerblank[18mm].
  \]
\end{practicebox}

\begin{practicebox}[Шаг 3. Вложенные отрезки]
  Порядок $M-P-Q$. $MP=3$ см, $MQ=8$ см.
  Отрезок $MP$ входит в \answerblank[18mm], поэтому
  \[
    PQ=\answerblank[18mm]-\answerblank[18mm]
      =\answerblank[18mm]\text{ см}.
  \]
\end{practicebox}

\begin{practicebox}[Шаг 4. Исправьте]
  Ученик написал: при порядке $K-M-N$, $KM=4$ см и $KN=10$ см,
  $MN=4+10=14$ см.

  Причина: \answerblank[95mm]

  Исправление: \answerblank[95mm]
\end{practicebox}

\begin{checkpointbox}[Повторная проверка]
  Порядок $A-D-C-B$, $AB=16$ см, $AD=5$ см, $CB=3$ см.
  Запишите равенство частей и найдите $DC$. \writinglines{3}
\end{checkpointbox}
